\subsection{Experiment 4E: control-relevant identification diagnosis}
\label{sec:experiment4e}
\paragraph{Question and controlled comparisons.}
Does redesign fail because of numerical solution error, nonlinear training
amplitudes, or limitations already present in pooled linear-plant identification?
The same seed-1 fleet, collection controller, two training profiles, and held-out
moderate S-curve are retained. Compare original nonlinear-data regression,
column-scaled SVD solving the same ridge objective, SVD using only near-linear
nonlinear trajectories, and SVD using linear-plant trajectories at all three
amplitudes. The last case uses exact frozen-speed ZOH propagation for training.
All redesigned controllers are tested on the nonlinear plant.

For regression matrix $X$, define $S_{jj}=\|X_{:j}\|_2$. The numerical check is
\begin{equation}
 \min_W\left\|\begin{bmatrix}XS^{-1}\\\sqrt{\epsilon}S^{-1}\end{bmatrix}W
 -\begin{bmatrix}Y\\0\end{bmatrix}\right\|_F^2,
 \qquad \widehat\Theta^{\mathsf T}=S^{-1}W,\quad\epsilon=10^{-10}.
\end{equation}
Scaling changes the solver, not the regularization objective. M1, M2, and M5
retain their original solver. A small absolute ridge coefficient is not
necessarily negligible relative to weak regressor directions.

\paragraph{Results.}
\Cref{tab:experiment4e} shows that SVD reproduces the original M4 failure.
Pooled raw condition numbers are approximately $1.46$--$1.80\times10^6$,
falling to $5.86$--$7.76\times10^3$ after column scaling. All matrices have
full column rank, but weak directions remain. The solver change alters M4
tracking by less than $10^{-9}\,\mathrm{deg\,s^{-1}}$. Numerical roundoff
therefore does not explain the observed failure.

\begin{table}[htbp]
\centering
\caption{Experiment 4E: M4 diagnosis. Tracking is moderate-test yaw RMSE in
$\mathrm{deg\,s^{-1}}$. Relative $B$ error is averaged over clients and 161
speeds against each client's true small-signal matrix.}
\label{tab:experiment4e}
\begin{tabular}{lrrr}
\toprule
Identification data and solver & Tracking & Worst $\rho$ & Relative $B$ error\\
\midrule
Nonlinear, all amplitudes, original & 0.05737 & 1.0321 & 0.4633\\
Nonlinear, all amplitudes, SVD & 0.05737 & 1.0321 & 0.4633\\
Nonlinear, near-linear only, SVD & 0.06255 & 0.9770 & 0.4715\\
Linear plant, all amplitudes, SVD & 0.06147 & 0.9772 & 0.4679\\
\bottomrule
\end{tabular}
\end{table}

The original worst M4 instability occurs at $12\,\mathrm{m\,s^{-1}}$ for
handling\_07. Near-linear training restores frozen M4 stability, but does not
improve tracking. Linear-plant training also leaves substantial input-matrix
error and poor tracking. Global LPV remains unstable in both ablations.
Nonlinear tire mismatch is therefore not the sole explanation. Relative matrix
errors use the fixed state coordinates and are not coordinate-invariant metrics.

\Cref{fig:experiment4e} relates input-channel error, gain, and stability for a
representative heavy client. Complete client and speed results accompany the
report as CSV files. Inspection confirms the same $r-r^\star$ integrator
convention, radian states and inputs, and prefilter equilibrium equation in
synthesis and simulation. Existing controller tests pass, providing no evidence
of a sign or unit error without claiming exhaustive software verification.

\begin{figure}[htbp]
\centering
\includegraphics[width=\linewidth]{experiment_4e_diagnosis.pdf}
\caption{Experiment 4E for heavy\_00. Original and SVD all-amplitude curves
overlap. This representative client is not the fleet's worst-case client.}
\label{fig:experiment4e}
\end{figure}

\paragraph{Interpretation and decision.}
Training uses two fixed speed--maneuver pairings, not their full cross product.
Pooled clients have different dynamics even within a family, and the reciprocal
basis only approximates exact discrete-time dynamics. Weak excitation, pooling
mismatch, basis approximation, and ridge bias can interact. Their individual
causal contributions have not all been isolated. The demonstrated finding is
that a solver change or reduced tire excitation alone does not recover the
desired performance. Experiment 4F reduces the degrees of freedom through
physical relationships while keeping controller weights fixed.

A further linear-data check fits nominal\_00, heavy\_00, and handling\_00
separately using their own trajectories. Mean relative $B$ errors are
$2.29$--$4.20\times10^{-5}$ without ridge and
$2.68$--$3.83\times10^{-4}$ with the original ridge. These are far below
the pooled M4 error of $0.468$. This strongly implicates heterogeneous pooling
under the existing closed-loop data distribution, rather than an inability of
the basis to represent each individual client. It does not isolate pooling from
all interactions with excitation. In closed loop, regression coefficients can
compensate along correlated state--input directions, so the pooled best predictor
need not equal an average of the individually identified physical models.

\subsection{Experiment 4F: structured local-dynamics identification}
\label{sec:experiment4f}
\paragraph{Method and information boundary.}
Fit near-linear nonlinear-plant trajectories using three positive ratios
\begin{equation}
 p=(a_f,a_r,h)=\left(C_f/m,\ C_r/m,\ m/I_z\right).
\end{equation}
Only the common geometry $l_f=1.2\,\mathrm m$, $l_r=1.6\,\mathrm m$ is known
to the estimator. True client mass, inertia, and stiffnesses are not supplied.
The continuous matrices satisfy
\begin{equation}
 A_c(v;p)=\begin{bmatrix}
 -(a_f+a_r)/v & (a_rl_r-a_fl_f)/v^2-1\\
 h(a_rl_r-a_fl_f) & -h(a_fl_f^2+a_rl_r^2)/v
 \end{bmatrix},\qquad
 B_c(v;p)=\begin{bmatrix}a_f/v\\ha_fl_f\end{bmatrix}.
\end{equation}
This avoids attempting to recover individually unidentifiable absolute physical
quantities. The estimator minimizes one-step RK4 residuals at $0.01$ s, scaled
by $0.05^\circ$ and $0.5^\circ\mathrm{/s}$ for sideslip and yaw rate.
Logarithmic parameters enforce positivity, with bounds $a_f,a_r\in[1,300]$
and $h\in[0.05,3]$ in corresponding SI units. Two common starting points,
$(40,40,0.5)$ and $(80,60,1)$, are used for each group. Training loss selects
the solution, with no held-out tracking-based selection. These are effective
group fits, not claimed exact family-center recovery.

\paragraph{Architecture comparison and changed estimation procedure.}
One structured fit is obtained globally and one per oracle family. M3 and M4
evaluate them continuously in speed. M1 and M2 freeze the corresponding fit at
the preselected $20\,\mathrm{m\,s^{-1}}$. M5 uses five frozen family fits
with nearest-anchor switching. All architectures receive the same structural
prior. This differs from 4D's independently optimized constant and gridded
regressions: 4F isolates deployment architecture after common identification.
It does not establish dominance over every LTI fitting procedure. Earlier
4D and 4E baselines remain part of the record.

Synthesis uses exact frozen-speed ZOH matrices, unchanged $Q,R$, and the existing
gain grid. Tests verify the reduced matrices against the validated physics and
RK4 against ZOH. Scaled local residual-Jacobian condition numbers are
$1.33$--$1.87$. Both starts agree within $9\times10^{-8}$ relative difference.
These support local numerical identifiability, not a global proof.

\paragraph{Results and gate.}
\Cref{tab:experiment4f} gives the moderate results. M4 reduces mean tracking
error by $8.7\%$ relative to structured M2 and $9.1\%$ relative to structured
M3. It approaches the physics-transfer value $0.03464$ without blending with
oracle matrices. Near-linear and strong M4 tracking errors are $0.01644$ and
$0.05678\,\mathrm{deg\,s^{-1}}$. All architectures are amplitude-feasible
at all three severities.

\begin{table}[htbp]
\centering
\caption{Experiment 4F: structured identification and redesign. Tracking and
worst-family values are yaw RMSE in $\mathrm{deg\,s^{-1}}$.}
\label{tab:experiment4f}
\begin{tabular}{lrrr}
\toprule
Architecture & Mean tracking & Worst family & Worst frozen $\rho$\\
\midrule
M1 & 0.04065 & 0.04340 & 0.9619\\
M2 & 0.03775 & 0.03831 & 0.9646\\
M3 & 0.03792 & 0.04453 & 0.9643\\
M4 & \textbf{0.03446} & \textbf{0.03636} & 0.9663\\
M5 & 0.04228 & 0.04270 & 0.9663\\
\bottomrule
\end{tabular}
\end{table}

The initial recovery gate passes: M4 is feasible, passes the frozen-grid audit,
and improves tracking by over $5\%$ against both structured alternatives.
This is one fleet with unseen trajectories on the same clients. Frozen
stability is not a time-varying LPV certificate. Next, validate independent
fleet seeds and additional speed--maneuver combinations, including the original
fitted LTI baselines. Only then introduce local identification, federation, and
learned clustering. Structured local-dynamics identification is a promising
repair, not evidence of universal M4 superiority.
